% !TeX root = ../main.tex
% Add the above to each chapter to make compiling the PDF easier in some editors.

\chapter{Introduction}\label{chapter:introduction}

Predicting the runtime of a database query before executing it is a fundamental problem in database systems research.
Performance predictions enable query scheduling, resource scaling, and query optimization~\parencite{t3}.
A query optimizer, for instance, compares alternative execution plans and selects the one it estimates to be fastest.
This makes the quality of runtime predictions directly relevant to the efficiency of query execution.

Traditional cost models, such as the one built into PostgreSQL, rely on hand-crafted formulas and statistical summaries of the stored data~\parencite{postgresql_docs_planner_stats}.
However, these models struggle with complex queries.

Learned cost models address this by using machine learning to capture system behavior from execution traces.
Recent work shows that they substantially outperform hand-crafted models~\parencite{zeroshot}.
Neural network approaches achieve high accuracy but suffer from high prediction latency~\parencite{stage}.

Rieger and Neumann address this trade-off with T3 (Tuple Time Tree)~\parencite{t3}, a compiled gradient-boosted decision tree model for the Umbra database system.
T3 achieves state-of-the-art (SOTA) accuracy while being orders of magnitude faster than neural network models.
It relies on two core ideas: decomposing query plans into pipelines and predicting the per-tuple (per-row) execution time for each pipeline independently.
Rieger and Neumann explicitly note that, while T3 is implemented for Umbra, its concepts are transferable to other database systems~\parencite{t3}.

PostgreSQL is one of the most widely deployed open-source relational database systems~\parencite{dbengines}.
It uses the Volcano iterator model~\parencite{graefe1994} rather than Umbra's compiled execution model~\parencite{umbra2020}, and exposes different plan information~\parencite{postgresql_docs_query_path}.
This raises the question of whether T3's pipeline-based approach can effectively transfer to PostgreSQL.

\paragraph{Research Question.}
\textit{How accurately can a decision tree model predict query runtimes on an iterator-based database system such as PostgreSQL?
Which of T3's key concepts -- pipeline-based representation and tuple-centric prediction targets -- transfer effectively to this setting?
How competitive is the result against existing learned cost models?}

\paragraph{}
To answer this, we adapt T3 for PostgreSQL using the public ZeroShot benchmark data~\parencite{zeroshot}.
We convert parsed PostgreSQL execution plans into T3's pipeline representation and construct a PostgreSQL-native feature vector.
We then train and evaluate a LightGBM model~\parencite{lightgbm} on a leave-one-database-out split to test generalization across unseen database instances.

Our results show that the adapted model generalizes well within the ZeroShot benchmark workload.
Furthermore, on the Join Order Benchmark, our model outperforms an earlier T3-based PostgreSQL implementation by Wehrstein et al.~\parencite{jobcomplex}, indicating the relevance of a PostgreSQL-native feature vector.

Compared to a broader set of learned cost models, our model ranks among the top performers, achieving a median Q-Error of 2.42.
However, it still falls behind state-of-the-art neural network models and T3 (Umbra).

A key finding is that PostgreSQL's iterator-based execution model makes exact pipeline timing extraction infeasible, which limits the benefit of T3's per-pipeline and per-tuple prediction strategies.
Per-query prediction works best in our PostgreSQL setting.


